\subsubsection{CLAHE}
The CLAHE augmentation technique has two hyperparameters: the clipping limit and the filter size. The clipping limit sets the maximum threshold for contrast limiting, while the filter size determines the size of the grid used for histogram equalization. To tune these hyperparameters, we used log-scale values ranging from 2 to 32. Extreme values were intentionally avoided, as explained in \ref{hptuning}. The hyperparameter tuning results for CLAHE are presented in table \ref{tab:clahe_hp}. 
\begin{table}[ht]
    \centering
    \caption{Results of Hyperparmeter Tuning for CLAHE.}
    \begin{tabular}{llll}
    \toprule
        Index & Clipping Limit & Filter Size & F1\_Score\_Top3 \\ \midrule 
        0 & 2 & 4 & 0.098 \\ 
        1 & 2 & 8 & 0.0965 \\ 
        2 & 2 & 16 & 0.0357 \\ 
        3 & 2 & 32 & 0.0607 \\ 
        4 & 4 & 4 & 0.1111 \\ 
        5 & 4 & 8 & 0.0862 \\ 
        6 & 4 & 16 & 0.1189 \\ 
        7 & 4 & 32 & 0.085 \\ 
        8 & 8 & 4 & 0.1113 \\ 
        9 & 8 & 8 & 0.0893 \\ 
        10 & 8 & 16 & 0.067 \\ 
        11 & 8 & 32 & 0.0543 \\
        12 & 16 & 4 & \textbf{0.1194} \\ 
        13 & 16 & 8 & 0.0765 \\ 
        14 & 16 & 16 & 0.0393 \\ 
        15 & 16 & 32 & 0.0282 \\ 
        16 & 32 & 4 & 0.0536 \\ 
        17 & 32 & 8 & 0.09 \\ 
        18 & 32 & 16 & 0.0601 \\ 
        19 & 32 & 32 & 0.0333 \\ 
    \bottomrule
    \end{tabular}
    \label{tab:clahe_hp}
\end{table}

\subsubsection{Color Jittering}
Color Jittering transformation has four hyperparameters: brightness, hue, contrast, and saturation. The values of these parameters determine the upper limit of the random magnitude applied to the image. For instance, if brightness has a value of 0.5, the brightness factor will be randomly chosen between 0 and 0.5. Hence, it is meaningful not to explore extreme values during hyperparameter tuning. In our experiment, we explored the maximum bound value for all parameters within the range of 0.2 to 0.8. The resulting hyperparameter tuning outcomes are presented in Table \ref{tab:colorjittering_hp}.

\begin{table}[ht]
\centering
\caption{Result of Hyperparameter Tuning for  Color Jittering.}
\begin{tabular}{lccccc}
\toprule
Index & Brightness & Hue & Contrast & Saturation & F1 Score \\
\midrule
1 & 0.57 & 0.45 & 0.56 & 0.78 & 0.0125 \\
2 & 0.74 & 0.73 & 0.6 & 0.4 & 0.0250 \\
3 & 0.41 & 0.28 & 0.55 & 0.47 & 0.0388 \\
4 & 0.55 & 0.2 & 0.23 & 0.64 & 0.0418 \\
5 & 0.67 & 0.26 & 0.65 & 0.27 & 0.0507 \\
6 & 0.77 & 0.77 & 0.24 & 0.36 & 0.0518 \\
7 & 0.48 & 0.55 & 0.34 & 0.64 & 0.0519 \\
8 & 0.8 & 0.42 & 0.54 & 0.65 & \textbf{0.0585} \\
\bottomrule
\end{tabular}
\label{tab:colorjittering_hp}
\end{table}

\subsubsection{GridDropout}
For Griddropout, we can adjust the number of grids and the size of each hole. The hole ratio refers to the size of the hole relative to a fixed grid cell. However, these two hyperparameters are correlated because a smaller number of holes can result in different interpretations of the same hole ratio. Therefore, we fixed the number of grid holes to 100 and only varied the hole ratio during hyperparameter tuning. The resulting hyperparameter tuning outcomes are presented in Table \ref{tab:griddropout_hp}.

\begin{table}[ht]
    \centering
    \caption{Result of Hyperparameter Tuning for GridDropout.}
    \begin{tabular}{ll}
    \toprule
        Hole ratio & F1-Score Top-3  \\ \midrule
        0.6 & 0.1021  \\ 
        0.5 & \textbf{0.1116}  \\ 
        0.2 & 0.08904  \\ 
        0.3 & 0.1013  \\ 
        0.4 & 0.0914  \\
        \bottomrule
    \end{tabular}
    \label{tab:griddropout_hp}
\end{table}

\subsubsection{Sharpen}
In sharpen we have to determin the alpha range and the lightness range. Alpha is the range to choose the visibility of the sharpened image. At 0, only the original image is visible, at 1.0 only its sharpened version is visible. Lightness is the range to choose the lightness of the sharpened image. We do grid search for same-size range with different of 0.2 from 0.2 to 1.0. The results are presented in table \ref{tab:sharpen_hp}
\begin{table}[!ht]
    \centering
    \caption{Result of Hyperparameter Tuning for Sharpen.}
    \begin{tabular}{lll}
    \toprule
        Alpha range & Lightness range & F1-Score Top 3 \\ \midrule
        (0.6, 0.8) & (0.6, 1.0)& \textbf{0.0908} \\ 
        (0.4, 0.6) & (0.2, 0.6)& 0.0965 \\ 
        (0.6, 0.8) & (0.4, 0.8)& 0.0075 \\ 
        (0.2, 0.4) & (0.6, 1.0) & 0.0774 \\ 
        (0.4, 0.6) & (0.4, 0.8)& 0.0724 \\ 
        (0.2, 0.4) & (0.4, 0.8)& 0.0878 \\ 
        (0.2, 0.4) & (0.4, 0.8) & 0.0624 \\ 
        (0.6, 0.8) & (0.2, 0.6) & 0.0568 \\ 
        (0.2, 0.4) & (0.2, 0.6) & 0.0849 \\ 
        (0.4, 0.6) & (0.6, 1.0) & 0.0477 \\ 
        \bottomrule
    \end{tabular}
    \label{tab:sharpen_hp}
\end{table}

\subsubsection{Motion Blur}
\begin{table}[!ht]
    \centering
    \begin{tabular}{|l|l|}
    \hline
        blur\_limit & F1 Score Top 3 \\ \hline
        71 & 0.2689 \\ \hline
        61 & 0.2825 \\ \hline
        51 & 0.2707 \\ \hline
        41 & 0.2774 \\ \hline
        31 & 0.2831 \\ \hline
        21 & 0.3048 \\ \hline
        11 & 0.2872 \\ \hline
    \end{tabular}
\end{table}